\section{راهنمای قالب \lr{(Template Guide)} — قابل حذف در نسخه‌ی نهایی}

\begin{stepbox}
\field{۱. جعبه‌ی گام \lr{(stepbox)}}{
همین جعبه‌ای که این متن داخل آن است. برای هر تمرین/زیرمسئله از دستور \texttt{\textbackslash field\{عنوان\}\{متن\}} داخل \texttt{stepbox} استفاده کنید.
}
\end{stepbox}

\vspace{0.5cm}
\noindent برای نکات جانبی، به‌جای نوشتن آن‌ها داخل \texttt{stepbox} اصلی، از سه جعبه‌ی معنایی زیر استفاده کنید تا گزارش خواناتر شود:

\begin{notebox}
این یک \textbf{نکته‌ی توضیحی} است؛ مثلاً یادآوری یک تعریف ریاضی یا اشاره به مرجع بیرونی.
\end{notebox}

\begin{tipbox}[پیشنهاد عملی]
این یک \textbf{پیشنهاد} است؛ مثلاً یک ترفند برای سریع‌تر اجرا شدن کد یا انتخاب بهتر هایپرپارامتر.
\end{tipbox}

\begin{warnbox}
این یک \textbf{هشدار} است؛ مثلاً نقطه‌ای که در پیاده‌سازی معمولاً اشتباه رخ می‌دهد (مثل عدم نرمال‌سازی ماتریس لاپلاسین).
\end{warnbox}

\subsection{نمایش کد}
برای اسنیپت‌های پایتون یا پرامپت‌های مدل مولد تصویر از \texttt{codebox} به همراه \texttt{lstlisting} استفاده کنید:

\begin{codebox}[spectral\_clustering.py]
\begin{lstlisting}[language=Python]
import numpy as np
from sklearn.cluster import KMeans

def spectral_clustering(A, k):
    D = np.diag(A.sum(axis=1))
    L = D - A
    eigvals, eigvecs = np.linalg.eigh(L)
    U = eigvecs[:, :k]
    labels = KMeans(n_clusters=k, n_init=10).fit_predict(U)
    return labels
\end{lstlisting}
\end{codebox}

\subsection{معادلات}
معادلات همان‌طور که قبلاً می‌نوشتید کار می‌کنند و نیازی به تغییر ندارند:
\[
L = D - A, \qquad Lx = \lambda x
\]

\subsection{گالری تصاویر کنار هم}
برای مقایسه‌ی دو تصویر (مثلاً قبل/بعد از یک مدل مولد تصویر) به‌جای دو \texttt{figure} جدا از این الگو استفاده کنید:

\begin{figure}[h!]
\centering
\begin{subfigure}{0.46\textwidth}
    \centering
    \fbox{\rule{0pt}{3cm}\rule{3cm}{0pt}} % جای‌گیر تصویر ۱ — با \includegraphics جایگزین کنید
    \caption{قبل}
\end{subfigure}
\hfill
\begin{subfigure}{0.46\textwidth}
    \centering
    \fbox{\rule{0pt}{3cm}\rule{3cm}{0pt}} % جای‌گیر تصویر ۲ — با \includegraphics جایگزین کنید
    \caption{بعد}
\end{subfigure}
\caption{نمونه‌ی نمایش دو تصویر کنار هم برای مقایسه}
\label{fig:before_after_demo}
\end{figure}

\begin{notebox}[نحوه‌ی استفاده در گزارش‌های بعدی]
کافی است در ابتدای \texttt{report.tex}، پنج دستور \texttt{\textbackslash reportWeek}، \texttt{\textbackslash reportTitleFa} و... را با عنوان تمرین جدید عوض کنید و فایل‌های \texttt{steps/stepXX.tex} جدید را \texttt{input} کنید؛ کل ظاهر و استایل بدون تغییر باقی می‌ماند.
\end{notebox}
